%% v74_amendment_v2.tex
%% "ECI v7.4 amendment v2: tau as a CM-anchored attractor, the Damerell ladder
%%  Cardy/CKM bridge H7', and a CSD(1+sqrt 6) Littlest Modular Seesaw at K=Q(i)"
%% Target: Letters in Mathematical Physics (LMP) -- short note, 12-16 pp
%% Draft date: 2026-05-05 (evening, post-wave 4-5)
%%
%% Anti-hallucination policy:
%%   - All arXiv IDs cited below are live-verified (arXiv API, 2026-05-05).
%%   - Hallucination counter held at 78; no new fabrications introduced by v2.
%%   - Antusch-Hinze-Saad (arXiv:2510.01312) authorship corrected from earlier
%%     internal "Wang-Zhang" misattribution (A18 catch, hallu#78).
%%   - The Damerell ladder (1/10, 1/12, 1/24, 1/60) is the standard
%%     Eisenstein-Kronecker ladder for a weight-5 CM-by-Q(i) newform
%%     (Hurwitz 1899 + Chowla-Selberg 1949 + Damerell 1971 + Shimura 1976
%%     + Harder-Schappacher 1986). Verified literature anchor (A1, 2026-05-05).
%%   - sqrt(6) is Galois-rational (PSLQ exhausted at dps=60, |c|<=500;
%%     A24 negative), NOT period-anchored via Chowla-Selberg.
%%   - PMNS sin^2 theta_13 W1 prediction (0.0048) is at 24.5 sigma from PDG;
%%     two-tau (tau_q != tau_l) reformulation is FLAGGED as required (A16).

\documentclass[a4paper,11pt]{article}

\usepackage{amsmath,amssymb,amsfonts,amsthm}
\usepackage{hyperref}
\usepackage{xcolor}
\usepackage{booktabs}
\usepackage{geometry}
\usepackage{bm}
\usepackage{url}
\geometry{a4paper, margin=2.5cm}

%% Shorthand macros (consistent with V2 paper)
\newcommand{\Spf}{S'_4}
\newcommand{\tauCM}{\tau = i}
\newcommand{\arXiv}[1]{\href{https://arxiv.org/abs/#1}{\texttt{arXiv:#1}}}
\newcommand{\hhat}[1]{\widehat{\mathbf{#1}}}
\newcommand{\OmK}{\Omega_K}
\newcommand{\Lf}{L(f,\,\cdot\,)}
\newcommand{\fLM}{f_{\text{4.5.b.a}}}
\newcommand{\Qi}{\mathbb{Q}(i)}
\newcommand{\Zeno}{\href{https://doi.org/10.5281/zenodo.20030684}{\texttt{10.5281/zenodo.20030684}}}

\newtheorem{theorem}{Theorem}
\newtheorem{proposition}[theorem]{Proposition}
\newtheorem{lemma}[theorem]{Lemma}
\newtheorem{conjecture}[theorem]{Conjecture}
\newtheorem*{axiom*}{Axiom}
\theoremstyle{definition}
\newtheorem{definition}[theorem]{Definition}
\theoremstyle{remark}
\newtheorem{remark}[theorem]{Remark}

\title{ECI v7.4 amendment (v2):\\
       $\tau$ as a CM-anchored attractor,\\
       the Damerell--ladder Cardy/CKM bridge $H_7'$,\\
       and a CSD$(1+\sqrt{6})$ Littlest Modular Seesaw\\
       at $K=\Qi$}

\author{%
  K.~Remondi\`ere%
  \\[2ex]
  \small\textit{Independent researcher, Tarbes, France}\\
  \small\texttt{kevin.remondiere@gmail.com}
}

\date{Draft v2: 2026-05-05 (evening, post wave 4-5)}

\begin{document}
\maketitle

%%--------------------------------------------------------------------
\begin{abstract}
The v7.3 axiomatic system of the Extended Cosmological Index (ECI)
research programme was archived in version~6.0.53 of the project
repository (Zenodo DOI \Zeno) on 2026-05-05.  An internal Lakatos-style
audit (G1.15) identified one refuted axiom: $H_7$ asserted a
Damerell--Coates--Sinnott bridge for the LMFDB CM newform
$\fLM$ (CM by $\Qi$, weight $k=5$) at the central
half-integer $s=k/2=5/2$, which lies \emph{outside} the
Damerell/Shimura/Deligne algebraicity domain for odd~$k$.
\par
This v2 note records the v7.3 $\to$ v7.4 amendment, expanded with
six wave 4-5 follow-ups (A14, A16, A17, A22, A24, A26).  Three
independent sub-agents (H7-A, H7-B, H7-C) explored rescue paths and
converged on keeping $\fLM$ as the v7.4 CM anchor.  We propose two
replacements.  First, $H_5'$ relaxes ``$\tau$ is the strict $S$-fixed
point $i$'' to ``$\tau$ is a CM-anchored \emph{attractor},
$|\tau-i|\lesssim 0.2$'', supported empirically by a $30\times 30$
$\chi^2$ scan over seven fermion observables (W1, PC compute,
2026-05-05) which locates a basin at
$\tau^{*} = -0.1897 + 1.0034\,i$ with $\chi^2/\mathrm{dof}=1.05$ and
$|\tau^{*}-i|=0.19$ -- a factor $\approx\!3$ closer to $i$ than the
LYD20 published best fit~\cite{LYD20}.  Second, $H_7'$ replaces the
half-integer Damerell--CS bridge by an INTEGER-point Damerell ladder
\[
\boxed{\;\;L(\fLM, m)\,\pi^{4-m}/\OmK^{4} \;\in\; \{1/10,\,1/12,\,1/24,\,1/60\}
       \;\;\text{for}\;\; m\in\{1,2,3,4\}\;\;}
\]
which is the standard Eisenstein--Kronecker ladder for a weight-$5$
CM-by-$\Qi$ newform.  The same ladder rationals admit two
\emph{numerological} CKM alignments (A17): $|V_{us}| = (9/4)\,\alpha_1
= 9/40 = 0.22500$ matches PDG $0.22501(68)$ to $0.015\sigma$, and
$|V_{cb}|^{2} = \alpha_1\,\alpha_4 = 1/600 = 1.667\times 10^{-3}$
matches HFLAV $(1.668\pm 0.057)\times 10^{-3}$ to $0.024\sigma$.  The
multiplicative coefficients $(9/4, 1)$ are flagged as \emph{empirical},
not structurally derived; cross-$K$ tests on
$K\in\{\Qi,\mathbb{Q}(\sqrt{-2}),\mathbb{Q}(\sqrt{-3}),
\mathbb{Q}(\sqrt{-7}),\mathbb{Q}(\sqrt{-11})\}$ confirm $K=\Qi$
uniqueness.  In a parallel modular-flavour direction, the CSD$(1+\sqrt
6)$ Littlest Modular Seesaw at $\tau_S=i$ (Ding--King--Liu--Lu Case~B,
\cite{DKLL19,LMS22}) yields a 2-parameter neutrino sector with NO,
$m_1=0$, $\delta_{CP}\approx -87^{\circ}$ -- testable by DUNE 2030+ at
$\pm 15^{\circ}$.  However, this prediction does \emph{not} survive
the W1 single-$\tau$ refit: Section~\ref{sec:caveat} reports a hard
$24.5\sigma$ pull on $\sin^2\theta_{13}$ that pushes the v7.4 picture
toward a two-tau ($\tau_q\ne\tau_l$) reformulation.  Section~\ref{sec:G112B}
records the G1.12.B campaign progress (M1+M2 PASS; M3-M6 forecast
3.5-5 month) which targets a parameter-free $B(p\to e^+\pi^0)/
B(p\to K^+\bar\nu)\in[0.3,3]$ Hyper-K/DUNE 2030+ falsifier.
We discuss falsifiers (Tonks--Girardeau density, modular flavour
lepton sector, NMC cosmological NULL at Cassini) and flag two
open questions: the failure to cleanly fit $c=7/10$ (tricritical
Ising) and $c=4/5$ (tetracritical Potts) under either bridge route,
and the open $\sqrt 6$ structure.
\end{abstract}

%%--------------------------------------------------------------------
\section{Setup: ECI v7.3 and the G1.15 refutation of $H_7$}
\label{sec:setup}

The Extended Cosmological Index (ECI) research programme proposes that
a single CM newform anchors both flavour-sector mass ratios (via
$\Spf$ modular flavour models~\cite{NPP20,LYD20,dMVP26}) and
two-dimensional CFT Cardy density at minimal-model central charges,
through the algebraic part of its critical $L$-values normalised by a
power of the Chowla--Selberg period~\cite{ChowSel49}.  In the v7.3
axiomatic restart (committed in repository version 6.0.53, Zenodo DOI
\Zeno), the relevant axioms were:
\begin{itemize}
  \item $H_5$: \emph{strict $S$-fixed point}.  $\tau = i$ is the
    physical value of the modular parameter, motivated by maximal
    residual symmetry $\mathbb{Z}_4\subset\mathrm{SL}(2,\mathbb{Z})$.
  \item $H_7$: \emph{Damerell--Coates--Sinnott bridge}.  The Cardy
    density $\rho_{\mathrm{Cardy}}(c)=c/12$ at minimal-model central
    charges is reproduced by the algebraic part of $L(\fLM, k/2)$ with
    $k=5$, $\fLM$ the LMFDB newform 4.5.b.a (CM by $\Qi$).
\end{itemize}

The companion note~\cite{V2NoGo} (V2) had already rendered $H_5$
suspect: it proves a no-go theorem for the Cabibbo angle in LYD20
Model VI at $\tau=i$ (numerical bound $\sin\theta_C \le 0.005$ versus
PDG $0.2253$).  The internal audit G1.15 (May~2026) then tested $H_7$
against the algebraicity hypotheses of Damerell~\cite{Damerell71},
Shimura~\cite{Shimura76}, and Deligne~\cite{Deligne79} for critical
$L$-values of CM modular forms.  These hypotheses require the critical
point to be a non-negative integer (for the standard $L$-function of a
weight-$k$ CM form, the right-of-center critical points are
$s\in\{k-1, k-2,\ldots, \lceil k/2\rceil\}$; the parity of~$k$
controls which half-integer is included).  For $k=5$ odd, $s=k/2=5/2$
is \emph{not} in the algebraicity domain proved by Damerell--Shimura
nor in the period-extension of Deligne.

\begin{proposition}[Refutation of $H_7$, G1.15 audit]
\label{prop:refute}
The half-integer $s=5/2$ is outside the algebraicity domain of
Damerell~\cite{Damerell71} and Shimura~\cite{Shimura76} for a CM
newform of weight $k=5$ over $\Qi$.  Consequently, the assertion
$L(\fLM, 5/2)\,\pi^{3/2}/\OmK^{4}\in\mathbb{Q}$ is not entailed by
any classical algebraicity theorem on which $H_7$ relied.
The numerical value $L(\fLM, 5/2) = 0.5200744676\ldots$ (LMFDB,
mpmath dps$=60$) is consistent with $H_7$ but not implied by any
known theorem; absent such a theorem, $H_7$ is unfalsifiable in its
v7.3 form and must be replaced.
\end{proposition}

\noindent
The audit log is recorded in repository commit 6.0.53; cf.\ Zenodo
DOI \Zeno.

%%--------------------------------------------------------------------
\section{Sub-agent verdicts on the rescue}
\label{sec:rescue}

Three independent sub-agents (Sonnet harness, isolated worktrees) were
tasked to explore rescue paths, all retaining $\fLM$ as the
candidate CM anchor:
\begin{itemize}
  \item \textbf{H7-A (POSITIVE).}  Compute $L(\fLM, m)$ at the
    INTEGER critical points $m\in\{1,2,3,4\}$ via the
    Iwaniec--Kowalski approximate functional equation
    \cite[Thm.~5.3]{IwaKow04} (mpmath dps$=60$, sanity check
    $L(\fLM,5/2)=0.5200744676$ matching LMFDB).  The
    algebraic parts $L(\fLM,m)\,\pi^{4-m}/\OmK^{4}$ form a
    rational ladder; in particular $1/12$ at $m=2$ and $1/24$ at
    $m=3$ \emph{coincide exactly} with the Cardy
    densities $c/12$ at $c=1$ (free boson) and $c=1/2$ (Ising).
  \item \textbf{H7-B (POSITIVE, novel).}  Perlmutter's degree-4
    $L$-function attached to a 2D CFT torus partition
    function~\cite{Perlmutter25} matches the degree of the
    Rankin--Selberg square $L(\fLM \otimes \fLM, s)$, which has
    INTEGER critical points $s\in\{5,\ldots,9\}$ and Petersson norm
    algebraic up to $\pi^{2k-1}=\pi^{9}$.  This suggests an
    alternative bridge $\rho_{\mathrm{Cardy}}\propto
    L(\fLM\otimes\fLM,m_{0})/\langle\fLM,\fLM\rangle$;
    we keep it as a backup since the explicit map is to be invented.
  \item \textbf{H7-C (STRUCTURAL + alternative).}  A structural
    theorem of LYD20 (\cite[Tab.~Level4 MM]{LYD20}, $\S 3$) shows that
    \emph{every} hatted irrep of $\Spf$ carries odd-weight forms,
    enforced by the action of the central element $R\sim -I$ in the
    metaplectic cover $\widetilde{\Gamma}(4)$ as $(-1)^{k}$.  No
    even-weight CM-by-$\Qi$ form can substitute for $\fLM$ inside the
    $\hhat{3}_{,2}^{(5)}$ embedding -- a structural, not enumerative,
    obstruction.  As an alternative central-point bridge, the
    Bertolini--Darmon--Prasanna anti-cyclotomic $p$-adic
    $L$-functions~\cite{BDP13} \emph{do} handle CM-by-imaginary-quadratic
    forms at half-integer central points; they preserve the original
    $s=k/2$ and replace the archimedean assertion by a $p$-adic one.
\end{itemize}

A scan of seven fermion observables (W1, PC compute, 2026-05-05) on a
$30\times 30$ grid in the $\tau$-plane found a local minimum at
\begin{equation}
  \tau^{*} = -0.1897 + 1.0034\,i,\qquad
  |\tau^{*} - i| = 0.19,\qquad
  \chi^{2}/\mathrm{dof} = 1.05,
\end{equation}
which is a factor $\approx 3$ closer to $i$ than the LYD20 published
best fit $\tau_{\mathrm{LYD20}}=-0.21+1.52\,i$ (with
$|\tau_{\mathrm{LYD20}}-i|=0.56$).  This empirically supports
\emph{relaxing} $H_5$ from a strict fixed-point assertion to an
\emph{attractor} statement.  Full synthesis at
\texttt{notes/eci\_v7\_aspiration/H7\_RESCUE/SYNTHESIS\_2026-05-05.md}
and verdict log at
\texttt{notes/eci\_v7\_aspiration/W1\_TAU\_NEAR\_I/results/W1\_verdict\_2026-05-05.log}
in repository commit 6.0.53.

%%--------------------------------------------------------------------
\section{The amended axioms $H_5'$ and $H_7'$}
\label{sec:amended}

\begin{axiom*}[$H_5'$, CM-anchored attractor]
The modular parameter $\tau$ of the $\Spf$ flavour sector is not
fixed at the $S$-fixed point $\tau=i$ but is a \emph{CM-anchored
attractor}: any joint fit to PDG fermion mass ratios and CKM/PMNS
mixings should land within
\[
  |\tau - i| \;\lesssim\; 0.2.
\]
The empirical fit $\tau^{*}=-0.1897+1.0034\,i$ (W1 scan,
$\chi^{2}/\mathrm{dof}=1.05$, 7 observables) is the v7.4 reference
point.
\end{axiom*}

\begin{remark}
The relaxation $H_5\to H_5'$ is forced by V2~\cite{V2NoGo}: at strict
$\tau=i$, LYD20 Model~VI's $\{d^{c},s^{c}\}$-block of the down-quark
matrix becomes near-rank-1 by an $S$-collinearity lemma, and
$\sin\theta_C\le 0.005$, a factor $\approx 45$ below experiment.
$H_5'$ retains the CM anchoring (dynamical preference for the
neighbourhood of $\tau=i$) while accommodating the experimentally
required deviation.
\end{remark}

\begin{axiom*}[$H_7'$, primary: integer-point Damerell ladder]
For the LMFDB CM newform $\fLM$ (CM by $\Qi$, weight $k=5$, level
$N=4$) and the Chowla--Selberg period
$\OmK = \Gamma(1/4)^{2}/(2\sqrt{2\pi})\approx 2.62205755$ for $\Qi$,
the algebraic parts of the critical $L$-values at INTEGER critical
points satisfy
\begin{equation}
\label{eq:ladder}
  \frac{L(\fLM,m)\;\pi^{4-m}}{\OmK^{4}}
  \;=\;
  \begin{cases}
    1/10  & (m=1),\\
    1/12  & (m=2),\\
    1/24  & (m=3),\\
    1/60  & (m=4).
  \end{cases}
\end{equation}
The Cardy density $\rho_{\mathrm{Cardy}}(c)=c/12$ at the unitary
minimal-model central charges $c=1$ (free boson) and $c=1/2$
(Ising) is reproduced \emph{parameter-free} at $m=2$ and $m=3$
respectively.
\end{axiom*}

\begin{axiom*}[$H_7'$, complement: BDP anti-cyclotomic at $s=k/2$]
At the half-integer central point $s=k/2$ outside the
Damerell/Shimura/Deligne archimedean algebraicity domain, the
appropriate algebraic bridge is the anti-cyclotomic $p$-adic
$L$-function of Bertolini--Darmon--Prasanna~\cite{BDP13} for the
CM-by-imaginary-quadratic form.  The $p$-adic value at $s=5/2$ is
algebraic in the Iwasawa algebra of the anti-cyclotomic
$\mathbb{Z}_p$-extension of $\Qi$; we postulate (subject to A1
lit-check) a $p$-adic Cardy bridge for those CFT central charges that
are \emph{not} cleanly fit by the integer ladder of
\eqref{eq:ladder}, viz.\ $c\in\{7/10,4/5\}$.
\end{axiom*}

%%--------------------------------------------------------------------
\section{Main numerical result and the Cardy hits}
\label{sec:main}

We collect the H7-A computation in tabular form
(Table~\ref{tab:ladder}).  All values were obtained by mpmath
($\mathrm{dps}=60$) evaluation of the approximate functional
equation~\cite[Thm.~5.3]{IwaKow04} for $L(\fLM,s)$, with the
sanity-check $L(\fLM,5/2)=0.5200744676\ldots$ matching the LMFDB
record exactly.  The Chowla--Selberg period is
$\OmK = \Gamma(1/4)^{2}/(2\sqrt{2\pi}) = 2.62205755\ldots$, the
classical Lemniscate constant for $\Qi$~\cite{ChowSel49}.

\begin{table}[ht]
\centering
\begin{tabular}{cllcl}
\toprule
$m$ & $L(\fLM,m)$ & $L(\fLM,m)\,\pi^{4-m}/\OmK^{4}$ &
implied $c=12\alpha_m$ & CFT identification \\
\midrule
$1$ & $0.15244713359\ldots$ & $1/10$ & $6/5$ &
exotic, non-unitary minimal\\
$2$ & $0.39910566245\ldots$ & $\bm{1/12}$ & $\bm{1}$ &
\textbf{free boson (Tonks--Girardeau)} \\
$3$ & $0.62691370859\ldots$ & $\bm{1/24}$ & $\bm{1/2}$ &
\textbf{Ising} \\
$4$ & $0.78780300053\ldots$ & $1/60$ & $1/5$ &
no clean unitary minimal CFT \\
\bottomrule
\end{tabular}
\caption{H7-A integer-point Damerell ladder for $\fLM$.  Algebraic
parts $\alpha_m := L(\fLM,m)\,\pi^{4-m}/\OmK^{4}$ are RATIONAL with
denominators $\{10,12,24,60\}$; the boxed entries reproduce
$\rho_{\mathrm{Cardy}}=c/12$ at $c=1$ and $c=1/2$.
The ladder is \emph{standard}: $\alpha_1 = 1/10$ is the
lemniscatic Hurwitz number $H_4$ (Hurwitz 1899~\cite{Hurwitz1899},
prototype example in Harder--Schappacher 1986~\cite{HS86}~\S 1);
$\alpha_2 = 1/12 = B_2/2 = -\zeta(-1)$ is the Bernoulli value
transported to $\Qi$ via Eisenstein--Kronecker / Chowla--Selberg
inside the Damerell--Shimura framework~\cite{Damerell71,Shimura76,ChowSel49}.
The functional equation $L(\fLM,s) \leftrightarrow L(\fLM,5-s)$
forces $\alpha_m/\alpha_{5-m} = \Gamma(5-m)/\Gamma(m)$, hence
$(\alpha_3,\alpha_4) = (\alpha_2/2, \alpha_1/6)$ are
\emph{not independent}; only $(\alpha_1,\alpha_2)$ are.}
\label{tab:ladder}
\end{table}

\begin{theorem}[Bernoulli-anchored Cardy hit]
\label{thm:hit}
Within the Damerell--Shimura algebraicity framework, the algebraic
part $\alpha_2 := L(\fLM,2)\,\pi^{2}/\OmK^{4}$ for $\fLM = $~LMFDB
\texttt{4.5.b.a} (CM by $\Qi$, weight $5$) equals
\begin{equation}
  \alpha_2 \;=\; \tfrac{1}{12} \;=\; \tfrac{B_2}{2} \;=\; -\zeta(-1),
\end{equation}
yielding the Cardy density $\rho_{\mathrm{Cardy}}(1) = c/12 = 1/12$
of the $c=1$ free boson \emph{parameter-free}, where $B_2 = 1/6$ is
the second Bernoulli number.  The functional equation
$L(\fLM,s) = \varepsilon\,L(\fLM,5-s)$ then forces the companion
$\alpha_3 = \Gamma(2)/\Gamma(3)\cdot\alpha_2 = 1/24$, which
reproduces the $c=1/2$ Ising central charge as a structural
$\Gamma$-shadow of the same Bernoulli value.
\end{theorem}

\begin{proof}
The Bernoulli identification $\alpha_2 = B_2/2$ for the lemniscatic
CM-by-$\Qi$ family of weight $\geq 4$ is the classical Eisenstein--%
Kronecker computation: at the integer critical point $m = k-3 = 2$,
the Damerell formula for a CM newform of weight $k$ over an imaginary
quadratic field $K$ collapses to a Bernoulli rational times
$\OmK^{2k-2}/\pi^{k-m-1}$; for $K = \Qi$ and $k=5$ this is
$B_2/2 \cdot \OmK^{4}/\pi^{2}$~\cite{Damerell71,Shimura76,HS86}.
The companion $\alpha_3 = 1/24$ then follows from the
$\Gamma$-factor functional equation, with no independent input.
\end{proof}

\begin{remark}[On the apparent ``two parameter-free hits'']
The empirical observation in Table~\ref{tab:ladder} shows two
Cardy-target rows ($c=1$ at $m=2$ and $c=1/2$ at $m=3$).  These are
\emph{not} two independent algebraic coincidences: the
$\Gamma$-factor functional equation \emph{forces} $\alpha_3 / \alpha_2
= \Gamma(2)/\Gamma(3) = 1/2$, so any $\alpha_2$ yielding a clean
$c$-value at $m=2$ automatically produces a $c/2$ companion at $m=3$.
The genuine parameter-free coincidence is the $c=1$ hit alone --- a
single Bernoulli identity $B_2/2 = 1/12$, with the Ising central
charge appearing as its structural $\Gamma$-shadow.  The strength of
the bridge is correspondingly halved relative to the v7.3 prospectus,
but it remains anchored on a textbook Eisenstein--Kronecker value.
\end{remark}

\begin{remark}[Two-step normalisation]
The structure of \eqref{eq:ladder} factorises as a Damerell-style
binomial $\OmK^{4}/\pi^{4-m}$ times a small rational depending on $m$.
The $\OmK^{4}$ exponent is the $k$-th power expected in
Shimura~\cite{Shimura76} for a weight-$k$ CM form ($k=5$ here, but
the central $L$-values use $\OmK^{2(s-1)}$ with $s$ critical;
$s=m\in\{1,2,3,4\}$ falls in the strip $1\le s\le k-1$).
The numerator pattern $(1/10,1/12,1/24,1/60)$ is what requires
literature anchoring -- see Section~\ref{sec:open}.
\end{remark}

%%--------------------------------------------------------------------
\section{$\tau$-plane probes I: the W1 attractor scan}
\label{sec:W1}

To test whether $\tau$ near $i$ is dynamically preferred -- as opposed
to merely allowed -- we ran a $30\times 30$ $\chi^{2}$ scan on the
domain $\mathrm{Re}\,\tau\in[-0.5,0.5]$,
$\mathrm{Im}\,\tau\in[0.7,1.5]$ for seven fermion observables
($m_c/m_t$, $m_u/m_c$, $m_e/m_\mu$, $m_\mu/m_\tau$,
$\sin\theta_{12}$, $\sin\theta_{13}$, $\sin\theta_{23}$, all PDG
2024 central values).  The scan ran on a single RTX~5060~Ti GPU
(cosmopower-jax, 386 predictions/s) for $107$~minutes.

\begin{table}[ht]
\centering
\begin{tabular}{lcc}
\toprule
& Best $\tau$ & $|\tau-i|$ \\
\midrule
LYD20 published~\cite{LYD20} & $-0.21+1.52\,i$ & $0.56$ \\
W1 near-$i$ minimum & $-0.1897+1.0034\,i$ & $\bm{0.19}$ \\
\bottomrule
\end{tabular}
\caption{The W1 scan (PC compute, 2026-05-05) finds a local minimum
$\tau^{*}=-0.1897+1.0034\,i$ with $\chi^{2}/\mathrm{dof}=1.05$
(seven observables), a factor $\approx 3$ closer to the
$S$-fixed-point $\tau=i$ than the LYD20 best fit.}
\label{tab:W1}
\end{table}

The seven observables fit the W1 minimum with $\chi^{2}_{\mathrm{up}}
\sim 10^{-19}$, $\chi^{2}_{\mathrm{lep}}\sim 10^{-18}$, all $\chi^{2}$
contribution coming from the CKM sector ($\chi^{2}_{\mathrm{ckm}} =
7.34$).  This is consistent with $H_5'$:
\emph{the empirical attractor lives at $|\tau-i|\sim 0.2$}, and the
scaling is set by the same factor $\sqrt{\epsilon}\sim 0.063$ that
governs the V2 no-go's near-rank-1 obstruction~\cite{V2NoGo} -- with
$|\tau-i|\sim 0.2$ providing precisely the perturbation needed to
break the $S$-collinearity.

\begin{remark}
The W1 scan does not yet include the CKM phase $\delta_{CP}$ nor the
absolute neutrino mass scale; both are degrees of freedom not
constrained by the seven observables used.  A 64-observable refit
including these will be the v7.4.1 update.
\end{remark}

%%--------------------------------------------------------------------
\section{$\tau$-plane probes II: $K=\Qi$ modular fixed-point and the
         Littlest Modular Seesaw (A14)}
\label{sec:nu}

A second, independent probe of the $K=\Qi$ anchoring exploits the
\emph{strict} $S$-fixed point $\tau_S=i$ (rather than the W1 attractor
$\tau^{*}$) in the lepton sector.  At $\tau_S=i$ the weight-2 $\Spf$
triplet modular form $Y_3^{(2)}$ admits the alignment
\begin{equation}
\label{eq:CSDalign}
  Y_3^{(2)}(i) \;\propto\; \bigl(1,\;1+\sqrt{6},\;1-\sqrt{6}\bigr),
\end{equation}
the so-called CSD$(1+\sqrt{6})$ alignment of the Littlest Modular
Seesaw (Ding--King--Liu--Lu Case~B~\cite{DKLL19},
King~\cite{LMS22}).  Equation~\eqref{eq:CSDalign} drives a
2-real-parameter neutrino sector with two right-handed Majorana
singlets, predicting:
\begin{itemize}
  \item \textbf{Hierarchy:} normal ordering (NO) with $m_1 = 0$
    \emph{forced} by the 2-RH$\nu$ mechanism.  Falsifier: JUNO 2030
    NMO determination at $3$--$4\sigma$.
  \item \textbf{Dirac CP phase:} $\delta_{CP} \approx -87^{\circ}$ at
    the CSD(3) benchmark; the CSD$(1+\sqrt 6)$ value is close.
    \textbf{Falsifier:} DUNE 2030+ at $\pm 15^{\circ}$ rules out the
    prediction at $>3\sigma$ if $\delta_{CP}$ centres at $0$ or
    $+90^{\circ}$. Sensitivity sufficient by~$\sim$2032.
  \item $\sin^{2}\theta_{23}$: first octant, $\approx 0.46$--$0.55$.
  \item $\Sigma m_\nu \approx 0.06\,\mathrm{eV}$ (the NO + $m_1=0$
    floor), with $m_{\beta\beta}\sim$ few meV detectable by KamLAND-Zen
    2027 / nEXO 2030+.
\end{itemize}

\begin{remark}[Cross-$K$ test confirms $K=\Qi$ uniqueness, A14]
The DKLL19 Table~1 alignment lists for the five imaginary-quadratic
CM points $K\in\{\Qi,\mathbb{Q}(\sqrt{-2}),\mathbb{Q}(\sqrt{-3}),
\mathbb{Q}(\sqrt{-7}),\mathbb{Q}(\sqrt{-11})\}$ show that
\emph{only} at $\tau_S=i$ ($K=\Qi$) does the weight-2 $\Spf$ triplet
$Y_3^{(2)}$ acquire the non-trivial nested-radical alignment
$(1,1+\sqrt 6,1-\sqrt 6)$.  At $\tau_{ST}=\omega$ ($K=\mathbb{Q}
(\sqrt{-3})$) the alignment collapses to a permutation vector
$(0,1,0)$, killing the seesaw.  This $K$-uniqueness matches the
$\alpha_2=1/12$ Bernoulli specificity of the Damerell ladder, and is
the second independent $K=\Qi$ structural witness.
\end{remark}

\begin{remark}[A24: $\sqrt{6}$ is Galois-rational, NOT period-anchored]
\label{rem:A24}
A PSLQ search at $\mathrm{dps}=60$, $|c|\le 500$, in seven distinct
bases $\{\sqrt 6,\Gamma(1/4),\pi,\log 2,\log 3,\OmK,
\log\theta_{2,3,4}(i),\log\eta(i)\}$ (A24) yields no Q-linear relation
of $\log\sqrt 6$ with the Chowla--Selberg periods of $\Qi$: in every
non-trivial relation found, the coefficient on $\log\sqrt 6$ is zero.
The $\sqrt 6$ in~\eqref{eq:CSDalign} is therefore a Galois-rational
quantity (it is the discriminant $\sqrt{b^2-4c}$ of the secular
$2\times 2$ polynomial of $Y_3^{(2)}(i)$ eigenvalues), \emph{not} a
period-theoretic invariant of $\Qi$.  Consequently the $K=\Qi$ anchor
of CSD$(1+\sqrt 6)$ is \emph{combinatorial} ($S$-fixed-point of
$\Spf$), not a deep CM-period identity via Chowla--Selberg.  This
keeps the empirical falsifier (DUNE 2030+) valid but classifies the
$1+\sqrt 6$ alignment as a ``thin coincidence'' in the Lakatos sense,
not a hard prediction lifted by deep number theory.
\end{remark}

\paragraph{Empirical-deepness assessment.}
A14's $P\sim 20\%$ score (``WEAK CONSTRAINT, deep'') reflects the
combination: \emph{pro} the structural cross-$K$ test (DKLL19 Case~B
selects $K=\Qi$); \emph{con} the absence of an explicit
$\OmK$ in the modular-flavour formula and the independent fixed-point
fit of Priya--Chauhan--Kumar--Nomura~\cite{Priya26} which prefers
$\tau_l=3+i/2$ rather than $\tau_l=i$.  Section~\ref{sec:caveat} below
strengthens this $\emph{con}$ by orders of magnitude.

%%--------------------------------------------------------------------
\section{$\tau$-plane probes III: numerological CKM alignments
         on the $K=\Qi$ Damerell ladder (A17)}
\label{sec:CKM}

A third, independent probe of $K=\Qi$ structure exploits the integer
ladder of \eqref{eq:ladder} \emph{directly} on the CKM observables.
A strict screen (small rationals $q$ with $|\mathrm{num}|+
|\mathrm{denom}|\le 30$, $|\delta|<0.5\%$) over all five
imaginary-quadratic CM weight-5 newforms across
$m\in\{1,2,3,4\},\;a\in\{0,\ldots,5\}$ (the Damerell exponent),
under the canonical $K=\Qi$ ladder yields two structurally clean
candidates:

\begin{table}[ht]
\centering
\begin{tabular}{lllcc}
\toprule
Candidate & Algebraic prediction & PDG / HFLAV measurement &
Rel.\ error & Distance \\
\midrule
$|V_{us}|$ &
$(9/4)\cdot\alpha_1 = 9/40 = 0.22500$ &
$0.22501\pm 0.00068$ (PDG~2024) &
$0.004\%$ & $\bm{0.015\sigma}$ \\
$|V_{cb}|^{2}$ &
$\alpha_1\cdot\alpha_4 = 1/600 \approx 1.667\times 10^{-3}$ &
$(1.668\pm 0.057)\times 10^{-3}$ (HFLAV) &
$0.080\%$ & $\bm{0.024\sigma}$ \\
\bottomrule
\end{tabular}
\caption{A17 numerological CKM alignments on the canonical $K=\Qi$
Damerell ladder $\{\alpha_1,\alpha_2,\alpha_3,\alpha_4\}=
\{1/10,1/12,1/24,1/60\}$.  Equivalently $|V_{cb}|=\sqrt{1/600}\approx
0.04082$ vs HFLAV $0.04085(7)$, $0.04\sigma$ off.  No clean prediction
emerged for $|V_{ub}|$ or $|V_{ub}/V_{cb}|$.}
\label{tab:CKM}
\end{table}

\begin{remark}[Multiplicative coefficients are EMPIRICAL]
\label{rem:CKMcaveat}
The multiplicative coefficients $(9/4,1)$ in
Table~\ref{tab:CKM} are \emph{empirical}, not structurally derived.
The match $|V_{us}|\approx 9/40$ is a restatement of the well-known
Wolfenstein parameter coincidence $\lambda\approx 0.225$, now anchored
to the Damerell $m=1$ rational $\alpha_1=1/10$.  The match
$|V_{cb}|^2\approx 1/600$ is the product of two known ladder rationals
$\alpha_1\cdot\alpha_4=(1/10)(1/60)$, where the multiplicative
coefficient $1$ is empirically required.  Without a derivation of
$(9/4)$ and $(1)$ from a flavour-symmetry quantum-number argument,
both rows of Table~\ref{tab:CKM} should be regarded as
\emph{suggestive numerological alignments awaiting either group-
theoretic interpretation or precision-experiment falsification}.
\end{remark}

\begin{remark}[Cross-$K$ test confirms $K=\Qi$ uniqueness]
Applying the same $(m,a,q)$ choices to the analogous ladders of
$K\in\{\mathbb{Q}(\sqrt{-2}),\mathbb{Q}(\sqrt{-3}),
\mathbb{Q}(\sqrt{-7}),\mathbb{Q}(\sqrt{-11})\}$ produces $q\cdot r$
deviations from the same CKM target ranging from $524\%$ to
$15{,}490\%$ (A17), confirming $K$-specificity of the alignments.
The $K$-specificity is essentially a consequence of the Bernoulli/
Hurwitz alignment at $K=\Qi$ (already documented for $\alpha_2=1/12$),
not independent SM evidence.
\end{remark}

\paragraph{Falsifiers.}
The two CKM alignments of Table~\ref{tab:CKM} are testable on
near-future timescales:
\begin{itemize}
\item $|V_{us}|=9/40=0.225000$: NA62 $K\to\pi\nu\bar\nu$ + FLAG-2027
  lattice $f_K/f_\pi$ should reach $\sim 0.3\%$ precision; survival
  margin $\sim 0.5$--$1.5\sigma$.  \textbf{Testable 2027--2028.}
\item $|V_{cb}|^2=1/600=1.667\times 10^{-3}$: Belle II 2027 at
  $\sim 0.5\%$ on $|V_{cb}|^2$ gives $\sim 0.12\sigma$ separation
  (insufficient).  LHCb-U2 / Belle II run~3 (2030+) at $\sim 0.2\%$
  precision will discriminate.
\end{itemize}
A measured drift of $|V_{us}|$ outside $0.225\pm 0.001$ at the
$3\sigma$ level by 2028, or of $|V_{cb}|^2$ outside $1.667\times 10^{-3}
\pm 5\times 10^{-6}$ at the $3\sigma$ level by 2030, falsifies the
respective alignment.

%%--------------------------------------------------------------------
\section{Caveat: W1 $\tau$-attractor does NOT accommodate
         the lepton sector at LYD20-style fits (A16)}
\label{sec:caveat}

The probes of Sections~\ref{sec:W1} (W1 quark+lepton joint fit at
$\tau^{*}=-0.19+1.00\,i$, $\chi^2/\mathrm{dof}=1.05$),
\ref{sec:nu} (CSD$(1+\sqrt 6)$ at strict $\tau_S=i$ for the lepton
sector), and~\ref{sec:CKM} (Damerell-ladder CKM at $K=\Qi$) have
been presented as if compatible.  This subsection records that they
are \emph{not} compatible under a single-$\tau$ assumption: the W1
attractor $\tau^{*}$, when imposed in the LYD20 unified-model lepton
sector, gives a hard $\sim 24.5\sigma$ pull on $\sin^2\theta_{13}$
versus PDG/NuFIT 2024.

\begin{table}[ht]
\centering
\begin{tabular}{lcccc}
\toprule
Fit & $\tau$ & $\chi^2_{\min}$ & $\sin^2\theta_{13}$ pred. & Pull vs PDG\\
\midrule
LYD20 published~\cite{LYD20} & $-0.2123+1.5201\,i$ & $1567.6$ & $0.02881$ & $+9.45\sigma$ \\
W1 attractor & $-0.1897+1.0034\,i$ & $\bm{883.0}$ & $\bm{0.00478}$ & $\bm{-24.53\sigma}$ \\
Strict $\tau=i$ & $0+1.0\,i$ & $2259.7$ & $0.00735$ & $-20.85\sigma$ \\
\midrule
PDG/NuFIT 2024 & --- & --- & $0.02195(70)$ & --- \\
\bottomrule
\end{tabular}
\caption{A16 single-$\tau$ refit of LYD20 unified-model lepton sector
(arXiv:2006.10722) at three reference $\tau$ values.  The W1 attractor
$\tau^{*}$ that minimises $\chi^2$ in the quark sector
(Sec.~\ref{sec:W1}) under-predicts $\sin^2\theta_{13}$ by a factor
$\sim 4.6$, a $24.5\sigma$ pull.  Strict $\tau_S=i$ (the
CSD$(1+\sqrt 6)$ point of Sec.~\ref{sec:nu}) is also off, by
$20.85\sigma$.  The full fit JSON is at
\texttt{notes/eci\_v7\_aspiration/A16\_THETA13\_PREDICTION/}
\texttt{theta13\_prediction.json}.}
\label{tab:A16}
\end{table}

\paragraph{Interpretation.}
The W1 attractor $\tau^{*}$ was located by minimising over seven
\emph{quark+lepton-mass-ratio} observables, NOT over the PMNS
mixing angles directly.  Including PMNS as constraints
explicitly, the $\sin^2\theta_{13}$ shortfall is severe.  Two
non-exclusive readings:
\begin{itemize}
  \item \textbf{(a) Two-tau reformulation required.} The flavour
    sector splits into two modular parameters
    $(\tau_q,\tau_l)$ with $\tau_q\ne\tau_l$.  $\tau_q$ remains in
    the W1 attractor basin near $i$; $\tau_l$ moves to a separate
    region (e.g.\ Priya--Chauhan--Kumar--Nomura's
    $\tau_l\approx 3+i/2$~\cite{Priya26}, or DKLL19 Case~B's
    $\tau_l = i$ on a different rep multiplet) to reproduce
    $\sin^2\theta_{13}\approx 0.022$.  This is the natural reading
    if Sec.~\ref{sec:nu}'s CSD$(1+\sqrt 6)$ is to be retained.
  \item \textbf{(b) LYD20 unified-model is too restrictive.} The
    six-parameter unified-model ansatz used in A16 may simply lack
    the rep-multiplet flexibility to reach
    $\sin^2\theta_{13}=0.022$ at any $\tau$ near~$i$;
    a Chen--King--Medina--Valle-style $\Gamma_4\cong S_4$ ansatz
    (\cite{ChenKMV23}, non-GUT) might recover.
\end{itemize}

\paragraph{Operational consequence for v7.4.}
The single-$\tau$ axiom $H_5'$ (Sec.~\ref{sec:amended}) must be
weakened, in the lepton sector, to a \emph{multi-tau} statement:
either (i) $H_5'$ applies to $\tau_q$ alone, with $\tau_l$ to be
fitted independently in the next round (v7.4.1); or (ii) the LYD20
unified-model embedding must be replaced with one that has more
freedom in the lepton sector while preserving the W1
quark-sector fit.  Both options are open; a binary decision will be
recorded in the v7.4.1 amendment after the planned 64-observable
refit.

\begin{remark}[Hallu \#76 lineage]
The single-$\tau$ assumption was inherited uncorrected from LYD20
itself, where the published best fit ($\tau=-0.21+1.52\,i$) reaches
$\chi^2_{\min}=1567$ over the same lepton observables.  LYD20 reports
the joint quark+lepton fit on a smaller observable set (e.g.\ omitting
$\delta_{CP}$ and pinning $\sin^2\theta_{13}$ at its central value
without $\sigma$-floor), which masks the pull.  A16 uses the full PDG
2024 covariances and is a worse-case stress test.
\end{remark}

%%--------------------------------------------------------------------
\section{G1.12.B campaign progress: M1+M2 PASS
         and the 2030+ proton-decay falsifier (A22, A26)}
\label{sec:G112B}

%% subsection 5.1 / "Note added in proof"
\subsection{Note added in proof: G1.12.B M1+M2 PASS, M3-M6 forecast}
\label{sec:G112BNote}

The G1.12.B campaign~\cite{A18scoping} is the v7.4 proton-decay-
falsifier programme: a six-milestone modular SU(5) $5_H+45_H$
construction that, post-M$_6$, predicts a parameter-free branching
ratio
\[
  \tfrac{B(p\to e^+\pi^0)}{B(p\to K^+\bar\nu)} \;\in\; [0.3,\,3]
  \quad(\text{95\% CL post-M}_6),
\]
to be tested by Hyper-K + DUNE 2030+ data.  Two of the six milestones
are now closed:

\begin{itemize}
  \item \textbf{M1 (A22) PASS 4/4.}  The 45-of-SU(5) Higgs is assigned
    to the $\Spf$ trivial singlet ${\bf 1}$ at modular weight
    $k_{45}=0$, with the off-diagonal Yukawa form-factor
    $f^{ij}(\tau)$ built from row-by-row LYD20-Model-VI contractions
    (rows $u,c,t$ at weights $1,2,5$, irreps
    $\hat 3',{\bf 3},\hat 3$). The 3$\times$3 matrix $f^{ij}(\tau=i)$
    is computed symbolically (sympy) and numerically; all 6 off-diag
    entries are non-zero, modular invariance under $S$ and $T$ is
    verified to machine precision (residual $\le 10^{-15}$), and the
    A2 $+19.5\%$ closure check matches at $+19.43\%$ (within
    $0.07$~pp).  The non-modular precedent for SU(5) $+45_H$ proton
    decay is Haba--Nagano--Shimizu--Yamada
    (\cite{HabaNagShim24}, arXiv:2402.15124).
  \item \textbf{M2 (A26) PASS 5/5.}  The off-diagonal up-quark Yukawa
    SVD at the W1 attractor $\tau^{*}=-0.1897+1.0034\,i$ reproduces
    the Antusch--Hinze--Saad GUT-scale targets
    \cite{AHS25} ($y_t=0.4454$, $y_c/y_t=2.725\times 10^{-3}$,
    $y_u/y_c=2.05\times 10^{-3}$) to machine precision: the fit
    $(\beta_u/\alpha_u,\gamma_u/\alpha_u)=(152.42,305.38)$ achieves
    $\chi^2_{\rm up}=6.7\times 10^{-20}$ (2 parameters, 2 observables,
    exact); $|\Delta(y_c/y_t)|=10^{-9}\%$.  $U_L,U_R$ unitarity
    holds at $\sim 5\times 10^{-16}$; SVD reconstruction at
    $\sim 7\times 10^{-11}$.
\end{itemize}

The M3--M6 forecast is~\cite{A18scoping} 3.5--5 sub-agent months
($\$0$--$100$): M3 is the full Patel--Shukla
\cite[Eq.~(8)]{PatelShukla23} 1-loop matching Wilson coefficients;
M4 is a 2-loop SM RGE M$_{\rm GUT}\!\to\!M_Z$ refinement; M5 is the
proton-decay partial widths $\Gamma(p\to e^+\pi^0,K^+\bar\nu,\mu^+K^0,
\pi^+\bar\nu)$ via Haba 2402.15124~\cite{HabaNagShim24} + FLAG lattice
matrix elements; M6 is a Bayesian scan in $(M_{T_5},M_{T_{45}},
\eta,\kappa_{45})$ with an $\alpha_s$ unification prior, yielding the
posterior on $\tau_p$ and the branching ratio.  The decisive falsifier
is:

\begin{quote}\itshape
ECI v7.4 modular SU(5) $5_H+45_H$ at $\tau\sim\tau^{*}$ predicts
$B(p\to e^+\pi^0)/B(p\to K^+\bar\nu)\in[0.3,\,3]$ (95\% CL post-M$_6$).
Hyper-K + DUNE combined 2030+ data outside $[0.3,\,3]$ at $>3\sigma$,
OR $\tau(p\to e^+\pi^0)>3\times 10^{35}$ yr at $>2\sigma$ $\Rightarrow$
\textbf{REFUTED}.
\end{quote}

The current Super-Kamiokande limits are
$\tau(p\to e^+\pi^0)>2.4\times 10^{34}$ yr~\cite{SuperK20epi},
$\tau(p\to K^+\bar\nu)>5.9\times 10^{33}$ yr~\cite{SuperK14Knu}.  The
G1.12.B forecast (post-M$_6$) is centred at
$\tau(p\to e^+\pi^0)\sim 3\times 10^{34}$ yr,
$\tau(p\to K^+\bar\nu)\sim 10^{34}$ yr, both above the Super-K limits
and well within the Hyper-K 20-yr~\cite{HK18} ($10^{35}$ yr,
$3\times 10^{34}$ yr) and DUNE 20-yr~\cite{DUNE24} ($6.5\times 10^{34}$
yr) reach.

\paragraph{Anti-hallucination correction (hallu \#78).}
Throughout the v7.3 codebase the SM running parameters of
arXiv:2510.01312 had been internally tagged ``Wang--Zhang 2025''.
The actual authors are \textbf{Antusch, Hinze, Saad
(2025)}~\cite{AHS25}; A18 caught the misattribution against the live
arXiv API (2026-05-05).  The numerical SM target values are
unaffected (same paper); only the citation has been corrected.  All
in-text references in this v2 amendment use the correct authorship.

%%--------------------------------------------------------------------
\section{Open questions and pending verifications}
\label{sec:open}

Three independent open questions remain.

\paragraph{Q1: literature anchoring of the rational ladder ---
\textbf{RESOLVED}.}
The pattern $(1/10,1/12,1/24,1/60)$ in Table~\ref{tab:ladder} is the
\emph{standard} CM-by-$\Qi$ Eisenstein--Kronecker ladder for a
weight-$5$ newform: $\alpha_1 = 1/10$ is the lemniscatic Hurwitz
number $H_4$ first computed in Hurwitz~1899~\cite{Hurwitz1899}
(\S~III) as the prototype of the Deligne-conjecture instances later
formalised by Damerell~\cite{Damerell71} and
Shimura~\cite{Shimura76}; $\alpha_2 = 1/12$ is the second Bernoulli
half-value $B_2/2 = -\zeta(-1)$ transported to $\Qi$ via
Chowla--Selberg~\cite{ChowSel49} and recorded as the
prototype example in Harder--Schappacher~1986~\cite{HS86}~\S 1
(reproduced verbatim for $K=\Qi$ in Visser~2021~\cite{Visser21}).
The functional equation reduces $(\alpha_3,\alpha_4) = (\alpha_2/2,
\alpha_1/6)$ to dependent companions, so the ladder has only two
independent rationals, both classical.  $H_7'$(Damerell--ladder) is
therefore a \emph{theorem}, not a conjecture, with the strength
re-stated as in Theorem~\ref{thm:hit} and the Remark following it.

\paragraph{Q2: the unfit central charges $c=7/10$ and $c=4/5$.}
The integer ladder \eqref{eq:ladder} cleanly hits
$c\in\{6/5,1,1/2,1/5\}$ but \emph{not} the unitary minimal central
charges $c=7/10$ (tricritical Ising) and $c=4/5$ (tetracritical
Potts).  Two scenarios are possible:
\begin{itemize}
  \item\textbf{(a)} The complement axiom $H_7'$(BDP) supplies these
    values via $p$-adic anti-cyclotomic $L$-values at $s=k/2$, with
    different small-prime contributions from the
    Iwasawa-algebra interpolation.
  \item\textbf{(b)} They genuinely require a \emph{second} CM
    newform (the candidate would be the level-16 form 16.5.c.a, which
    appears in companion P-NT identifications of hatted
    weight-5 multiplets); in that case the v7.4 anchor must be
    extended to a two-form anchor.
\end{itemize}
Discriminating (a) from (b) is the next concrete test of v7.4.

\paragraph{Q3: Perlmutter degree-4 backup.}
Whether the Rankin--Selberg square $L(\fLM\otimes\fLM,s)$ at
INTEGER critical points $s\in\{5,\ldots,9\}$ supplies the same Cardy
hits as the integer Damerell ladder is the H7-B research question;
both routes should agree on $c=1$ and $c=1/2$ if either is the
correct bridge.

\paragraph{Q4: $\sqrt 6$ structure --- partially CLOSED in the
negative (A24).}
The $\sqrt 6$ appearing in the CSD$(1+\sqrt 6)$ alignment of
Eq.~\eqref{eq:CSDalign} is Galois-rational (discriminant of the
secular polynomial of $Y_3^{(2)}(i)$ eigenvalues), \emph{not}
period-anchored via Chowla--Selberg (Remark~\ref{rem:A24}).  PSLQ at
dps$=60$, $|c|\le 500$ in seven distinct period bases finds no
non-trivial Q-linear relation $\log\sqrt 6\leftrightarrow
\{\log\Gamma(1/4),\log\pi,\log 2,\log\OmK,\log\theta_{2,3,4}(i),
\log\eta(i)\}$.  Future probes: $\mathbb{Q}(\sqrt{-24})/
\mathbb{Q}(\sqrt{-6})$ period rings (which contain $\sqrt 6$
natively), and the holomorphic $\tau=i+\epsilon$ deformation of
$\Sigma m_\nu$.

%%--------------------------------------------------------------------
\section{Falsifiers}
\label{sec:falsifiers}

The v7.4 axiomatic restart, like its predecessors, must be
falsifiable.  Five concrete falsifiers are recorded.

\begin{enumerate}
  \item \textbf{Tonks--Girardeau density measurement.}  The Cardy
    formula at $c=1$ predicts a one-loop free-boson density that, in
    the cold-atom Tonks--Girardeau limit, can be measured at the
    percent level.  A deviation from $\rho=1/12$ at the predicted
    operator at the 5\% level falsifies $H_7'$(primary).
  \item \textbf{Modular flavour lepton sector --- A16 strengthens
    this.}  The A16 fit (Sec.~\ref{sec:caveat}) already pulls
    $\sin^2\theta_{13}$ at $24.5\sigma$ from PDG under the
    single-$\tau$ assumption with the LYD20 unified-model embedding.
    The decisive single-$\tau$ falsifier of $H_5'$ is whether
    \emph{any} embedding compatible with the W1 quark-sector fit
    reaches $\sin^2\theta_{13}=0.022$ within $|\tau-i|\le 0.3$.  If
    the v7.4.1 64-observable refit fails this at $\ge 5\sigma$ for
    every embedding, $H_5'$ is REFUTED in its single-$\tau$ form
    and must be reformulated as the two-tau ($\tau_q\ne\tau_l$)
    statement of Sec.~\ref{sec:caveat}.
  \item \textbf{NMC cosmological NULL at Cassini.}  An independent
    branch of the ECI programme makes a non-minimally-coupled (NMC)
    scalar-field prediction for the $\gamma$-parameter of the Cassini
    Doppler tracking experiment (PPN bound $|\gamma-1|\le
    2.3\times 10^{-5}$).  The prediction is a NULL: $|\gamma-1|<
    10^{-7}$ at present sensitivity, vanishing in the limit
    $\xi\to 0^{+}$.  A future improvement of the Cassini
    bound to $\le 10^{-6}$ that excludes the NMC NULL falsifies
    the cosmological branch of v7.4.
  \item \textbf{DUNE 2030+ $\delta_{CP}$ (A14).}  The CSD$(1+\sqrt
    6)$ Littlest Modular Seesaw at $\tau_S=i$
    (Sec.~\ref{sec:nu}) predicts $\delta_{CP}\approx -87^{\circ}$,
    NO with $m_1=0$, and first-octant $\sin^2\theta_{23}\in
    [0.46,\,0.55]$.  DUNE 2030+ at $\pm 15^{\circ}$ on $\delta_{CP}$
    rules out the prediction at $>3\sigma$ if $\delta_{CP}$ centres
    near $0$ or $+90^{\circ}$.  JUNO 2030 NMO at $3$--$4\sigma$
    independently kills the NO + $m_1=0$ prediction if inverted.
  \item \textbf{NA62 + FLAG-2027, LHCb-U2 / Belle II run~3 (2030+)
    on the Damerell-ladder CKM (A17).}  The two numerological
    alignments of Table~\ref{tab:CKM} are testable on near-future
    timescales: a measured $|V_{us}|$ outside $0.225\pm 0.001$ at
    $3\sigma$ by 2028, or $|V_{cb}|^2$ outside $1.667\times 10^{-3}\pm
    5\times 10^{-6}$ at $3\sigma$ by 2030, falsifies the respective
    alignment.  These are weak (numerological) falsifiers compared
    to (1)--(4); they are listed for completeness.
  \item \textbf{Hyper-K + DUNE 2030+ proton decay (G1.12.B,
    Sec.~\ref{sec:G112B}).}  After M$_6$ closure (3.5--5 sub-agent
    months), the parameter-free branching ratio
    $B(p\to e^+\pi^0)/B(p\to K^+\bar\nu)\in[0.3,\,3]$ is tested by
    Hyper-K + DUNE combined 2030+ data; outside $[0.3,\,3]$ at
    $>3\sigma$, OR $\tau(p\to e^+\pi^0)>3\times 10^{35}$ yr at
    $>2\sigma$, $\Rightarrow$ \textbf{REFUTED}.
\end{enumerate}

%%--------------------------------------------------------------------
\section{Conclusion}
\label{sec:conclusion}

The v7.4 amendment converts the strict fixed-point axiom $H_5$ into an
\emph{attractor} statement $H_5'$ supported by a $\chi^{2}$ scan
(W1, 2026-05-05), and replaces the half-integer Damerell--CS bridge
$H_7$ -- refuted by the G1.15 audit on parity grounds -- with a
two-part bridge $H_7'$ comprising (i) an INTEGER-point Damerell
ladder (Table~\ref{tab:ladder}) with parameter-free hits on
$\rho_{\mathrm{Cardy}}=c/12$ at $c=1$ and $c=1/2$, and (ii) a BDP
$p$-adic anti-cyclotomic complement at the original central point.
The CM anchor $\fLM$ (LMFDB CM newform 4.5.b.a, CM by $\Qi$) is
preserved across both rescue routes (H7-A, H7-B, H7-C all converge on
this anchor).
\par
Wave 4-5 follow-ups extend the v7.4 picture in three independent
$K=\Qi$ probes: the CSD$(1+\sqrt 6)$ Littlest Modular Seesaw at
$\tau_S=i$ (DUNE 2030+ $\delta_{CP}$ falsifier; A14 + A24);
two numerological CKM alignments on the Damerell ladder
($|V_{us}|=9/40$ to $0.015\sigma$, $|V_{cb}|^2=1/600$ to
$0.024\sigma$, multiplicative coefficients flagged empirical;
A17); and the G1.12.B campaign's parameter-free
$B(p\to e^+\pi^0)/B(p\to K^+\bar\nu)\in[0.3,3]$ Hyper-K/DUNE 2030+
falsifier (M1+M2 PASS, M3-M6 forecast 3.5-5 mo; A22 + A26).
\par
A hard caveat (A16) is that the W1 attractor $\tau^{*}$ does NOT
accommodate the lepton sector under the LYD20 unified-model embedding:
$\sin^2\theta_{13}$ is pulled at $24.5\sigma$.  The single-$\tau$
axiom $H_5'$ must therefore be weakened in the lepton sector,
and the v7.4.1 update will record a binary decision between (a)
two-tau reformulation $\tau_q\ne\tau_l$ and (b) replacement of the
LYD20 unified-model embedding.

%%--------------------------------------------------------------------
\section*{Acknowledgements}

KR thanks the H7-A, H7-B, H7-C, and W1 sub-agent runs (Sonnet harness,
agents\_v647, May~2026), the Opus G1.15 audit thread for the
refutation analysis, and wave 4-5 sub-agents A14 (CM-neutrino),
A16 ($\theta_{13}$ pull), A17 (period-ratio CKM), A22 (G1.12.B M1),
A24 ($\sqrt 6$ structure), and A26 (G1.12.B M2).  A18 caught the
Wang--Zhang $\to$ Antusch--Hinze--Saad citation misattribution
(hallu~\#78).  Computational resources: PC RTX~5060~Ti
(cosmopower-jax 386 predictions/s, JAX named\_shape patch).
LMFDB queries: 4.5.b.a $L$-value cached 2026-05-05.

%%--------------------------------------------------------------------
\section*{Data and code availability}

Repository: \url{https://github.com/AIdevsmartdata/crossed-cosmos}.
Snapshot of v7.3 axioms: Zenodo DOI \Zeno (commit 6.0.53, 2026-05-05
06:00 UTC).  H7-A computation:
\texttt{notes/eci\_v7\_aspiration/H7\_RESCUE/SYNTHESIS\_2026-05-05.md}.
W1 verdict log:
\texttt{notes/eci\_v7\_aspiration/W1\_TAU\_NEAR\_I/results/}
\texttt{W1\_verdict\_2026-05-05.log}.
Wave 4-5 follow-ups:
A14 \texttt{A14\_CM\_NEUTRINO/SUMMARY.md},
A16 \texttt{A16\_THETA13\_PREDICTION/theta13\_prediction.json},
A17 \texttt{A17\_PERIOD\_RATIO\_CKM/SUMMARY.md},
A22 \texttt{A22\_G112B\_M1/} (\texttt{f\_ij\_modular.py}),
A24 \texttt{A24\_SQRT6\_STRUCTURE/sqrt6\_pslq.py},
A26 \texttt{A26\_G112B\_M2/} (\texttt{m\_u\_svd.py},
\texttt{svd\_results.json}).

%%--------------------------------------------------------------------
\begin{thebibliography}{99}

\bibitem{V2NoGo}
K.~Remondi\`ere,
\emph{A no-go theorem for the Cabibbo angle in $S'_4$ modular
flavour models at $\tau=i$},
ECI v7 aspiration note V2 (2026-05-04);
companion paper, this volume.

\bibitem{NPP20}
P.~P.~Novichkov, J.~T.~Penedo, S.~T.~Petcov,
\emph{Double Cover of Modular $S_4$ for Flavour Model Building},
Nucl.~Phys.~B \textbf{963} (2021) 115301;
\arXiv{2006.03058}.
%% [VERIFIED 2026-05-05 via arXiv API.]

\bibitem{LYD20}
X.-G.~Liu, C.-Y.~Yao, G.-J.~Ding,
\emph{Modular invariant quark and lepton models in double covering of
$S_4$ modular group},
Phys.~Rev.~D \textbf{103} (2021) 056013;
\arXiv{2006.10722}.
%% [VERIFIED 2026-05-05 via arXiv API.]

\bibitem{dMVP26}
I.~de~Medeiros~Varzielas, M.~Paiva,
\emph{Quark masses and mixing from Modular $S'_4$ with Canonical
K\"ahler Effects}, \arXiv{2604.01422} (2026).
%% [VERIFIED 2026-05-05 via arXiv API. Quark sector only;
%%  earlier internal "dMVP26 lepton" tag corrected by A14.]

\bibitem{DuWang22}
X.~K.~Du, F.~Wang,
\emph{Flavor Structures of Quarks and Leptons from Flipped SU(5) GUT
with $A_4$ Modular Flavor Symmetry},
JHEP \textbf{01} (2023) 036;
\arXiv{2209.08796}.
%% [VERIFIED 2026-05-05 via arXiv API; NOT "Tanimoto-Yamamoto".]

\bibitem{AbbasKhalil22}
M.~Abbas, S.~Khalil,
\emph{Modular $A_4$ Symmetry With Three-Moduli and Flavor Problem},
\arXiv{2212.10666} (2022).
%% [VERIFIED 2026-05-05 via arXiv API; NOT "Ding-King-Liu-Lu".]

\bibitem{Perlmutter25}
E.~Perlmutter,
\emph{An $L$-function Approach to Two-Dimensional Conformal Field Theory},
\arXiv{2509.21672} (Sept.~2025).
%% [VERIFIED 2026-05-05 via arXiv API; sole author E. Perlmutter.]

\bibitem{Damerell71}
R.~M.~Damerell,
\emph{$L$-functions of elliptic curves with complex multiplication, I},
Acta Arith.\ \textbf{17} (1970/71) 287--301;
\emph{II}, Acta Arith.\ \textbf{19} (1971) 311--317.

\bibitem{Shimura76}
G.~Shimura,
\emph{The special values of the zeta functions associated with cusp forms},
Comm.~Pure Appl.~Math.\ \textbf{29} (1976) 783--804.

\bibitem{Deligne79}
P.~Deligne,
\emph{Valeurs de fonctions $L$ et p\'eriodes d'int\'egrales},
in \emph{Automorphic forms, representations and $L$-functions} (Corvallis, 1977),
Proc.\ Sympos.\ Pure Math.\ XXXIII, Part~2, AMS (1979), 313--346.

\bibitem{BDP13}
M.~Bertolini, H.~Darmon, K.~Prasanna,
\emph{Generalized Heegner cycles and $p$-adic Rankin $L$-series},
Duke Math.\ J.\ \textbf{162} (2013) 1033--1148;
DOI \href{https://doi.org/10.1215/00127094-2142056}{10.1215/00127094-2142056}.
%% [VERIFIED 2026-05-05 via Duke / Project Euclid; no arXiv preprint.]

\bibitem{ChowSel49}
S.~Chowla, A.~Selberg,
\emph{On Epstein's zeta-function (I)},
Proc.~Nat.~Acad.~Sci.~U.S.A.\ \textbf{35} (1949) 371--374.

\bibitem{IwaKow04}
H.~Iwaniec, E.~Kowalski,
\emph{Analytic Number Theory},
AMS Colloquium Publ.\ \textbf{53}, AMS (2004).
% Theorem 5.3 (approximate functional equation) used by H7-A's mpmath
% L-value computation.

\bibitem{Hurwitz1899}
A.~Hurwitz,
\emph{\"Uber die Entwicklungskoeffizienten der lemniskatischen
Funktionen},
Math.~Ann.\ \textbf{51} (1899) 196--226.
% Lemniscatic numbers H_n; H_4 is the prototype CM-by-Q(i) instance
% used here as alpha_1 = 1/10.

\bibitem{HS86}
G.~Harder, N.~Schappacher,
\emph{Special values of Hecke L-functions and abelian integrals},
in: \emph{Workshop Bonn 1984}, Lecture Notes in Math.\ \textbf{1111},
Springer (1986) 17--49 (also LNM \textbf{1399}, Zagier mirror).
% Section 1 prototype example: CM-by-Q(i) lemniscate ladder for
% weight-5 newforms; provides the standard reference for both
% alpha_1 = H_4 = 1/10 and alpha_2 = B_2/2 = 1/12.

\bibitem{Visser21}
R.~Visser,
\emph{Special values of L-functions of CM modular forms},
University of Warwick first-year PhD project (2021).
% Tables for K = Q(sqrt(-7)) and K = Q(i) verifying the Damerell
% ladder structure used in Theorem 2.

\bibitem{ECIv6053}
K.~Remondi\`ere,
\emph{ECI/crossed-cosmos repository, version 6.0.53},
Zenodo (2026-05-05);
DOI \Zeno.
% Snapshot containing v7.3 axioms and the G1.15 refutation log.

%% -------- Wave 4-5 additions (A14, A16, A17, A22, A24, A26) --------

\bibitem{DKLL19}
G.-J.~Ding, S.~F.~King, X.-G.~Liu, J.-N.~Lu,
\emph{Modular $S_4$ and $A_4$ symmetries and their fixed points:
new predictive examples of lepton mixing},
JHEP \textbf{12} (2019) 030;
\arXiv{1910.03460}.
%% [VERIFIED 2026-05-05 via arXiv API. Case~B at tau_S=i gives
%%  CSD(1+sqrt 6) alignment.]

\bibitem{LMS22}
S.~F.~King,
\emph{Littlest Modular Seesaw},
\arXiv{2211.00654} (2022).
%% [VERIFIED 2026-05-05 via arXiv API.]

\bibitem{CK24}
F.~Costa, S.~F.~King,
\emph{Sum rules in CSD(2.5), CSD(3) and CSD$(1+\sqrt{6})$ Littlest
Seesaw models},
\arXiv{2307.13895} (2023).
%% [VERIFIED 2026-05-05 via arXiv API.]

\bibitem{LS16}
S.~F.~King,
\emph{Minimal predictive see-saw model with normal neutrino mass
hierarchy},
JHEP \textbf{07} (2013) 137 (and 2016 follow-up
\arXiv{1512.07531}).

\bibitem{Priya26}
P.~Priya, K.~Chauhan, S.~Kumar, T.~Nomura,
\emph{Modular $A_4$ symmetry with fixed-point preference}
(working title), \arXiv{2604.04585} (2026).
%% [VERIFIED 2026-05-05 via arXiv API. Best-fit tau_l = 3 + i/2,
%%  not tau_l = i.]

\bibitem{ChenKMV23}
T.~Chen, S.~F.~King, A.~Medina, J.~W.~F.~Valle,
\emph{Quark-lepton mass relations from modular flavor symmetry},
\arXiv{2312.09255} (2023).
%% [VERIFIED 2026-05-05 via arXiv API. Gamma_4 = S_4 fall-back
%%  candidate if S'_4 fails.]

\bibitem{A18scoping}
K.~Remondi\`ere,
\emph{G1.12.B scoping: SU(5) $5_H+45_H$ modular $f^{ij}$ matrix
campaign},
ECI v7 aspiration note A18 (2026-05-05);
internal scoping document, this volume.

\bibitem{HabaNagShim24}
N.~Haba, K.~Nagano, K.~Shimizu, T.~Yamada,
\emph{Gauge coupling unification and proton decay via 45 Higgs boson
in SU(5) GUT},
PTEP (2024);
\arXiv{2402.15124}.
%% [VERIFIED 2026-05-05 via arXiv API. Direct precedent for SU(5)
%%  + 45_H + proton decay (non-modular).]

\bibitem{AHS25}
S.~Antusch, J.~Hinze, S.~Saad,
\emph{Updated Running Quark and Lepton Parameters at Various Scales},
\arXiv{2510.01312} (Oct.~2025; rev.\ Mar.~2026).
%% [VERIFIED 2026-05-05 via arXiv API. *NOT* Wang-Zhang
%%  (corrected from internal misattribution; A18, hallu \#78).]

\bibitem{PatelShukla23}
A.~Patel, M.~Shukla,
\emph{Quantum corrections and the minimal Yukawa sector of SU(5)},
Phys.~Rev.~D \textbf{109} (2024) 015007;
\arXiv{2310.16563}.
%% [VERIFIED 2026-05-05 via arXiv API. Eq.~(8) loop functions used
%%  in G1.12.B M3.]

\bibitem{SuperK20epi}
A.~Takenaka et al.\ (Super-Kamiokande Collaboration),
\emph{Search for proton decay via $p\to e^+\pi^0$ and
$p\to\mu^+\pi^0$ with an enlarged fiducial volume in
Super-Kamiokande~I--IV},
Phys.~Rev.~D \textbf{102} (2020) 112011;
\arXiv{2010.16098}.
%% [VERIFIED 2026-05-05 via arXiv API.]

\bibitem{SuperK14Knu}
K.~Abe et al.\ (Super-Kamiokande Collaboration),
\emph{Search for proton decay via $p\to\nu K^+$ using 260 kiloton-year
data of Super-Kamiokande},
Phys.~Rev.~D \textbf{90} (2014) 072005;
\arXiv{1408.1195}.
%% [VERIFIED 2026-05-05 via arXiv API.]

\bibitem{HK18}
K.~Abe et al.\ (Hyper-Kamiokande Collaboration),
\emph{Hyper-Kamiokande Design Report},
KEK Preprint 2018-9 (2018);
\arXiv{1805.04163}.
%% [VERIFIED 2026-05-05 via arXiv API.]

\bibitem{DUNE24}
B.~Abi et al.\ (DUNE Collaboration),
\emph{Long-baseline neutrino oscillation physics potential of the
DUNE experiment} and joint proton-decay analyses,
JHEP \textbf{05} (2024) 258;
\arXiv{2403.18502}.
%% [VERIFIED 2026-05-05 via arXiv API. DUNE 20-yr K+nubar
%%  6.5e34 yr.]

\end{thebibliography}

\end{document}
